\section{Outstanding tasks}\label{sec:outstanding}

\subsection{Obtain a sign from the actual remaining arithmetic}

The central unpaid step is a useful lower bound for
\(\cD_a(q_{\can}^\circ)\) under the actual event equations,
exceeding the complete error \(\eta_a\) in \eqref{eq:eta}.
The phase law is
\[
 u_n^{\can}=\ee^{\ii\Theta(T,a)}
            \ee^{-\ii(a/2+\pi T/8)\log n}.
\]
Every phase, weight, cutoff, carrier, and correction comes from the same
\((a,T)\). Multiplicativity alone, even supplemented by the growing
calibration of Corollary~\ref{cor:growing} and the one mirror inequality,
does not enforce the desired lower bound.
A further argument must use information absent from
Theorem~\ref{thm:joint}, or select a materially different measurement.
Merely evaluating the field more accurately does not supply its sign.

The next bounded analytical question should specify the new information,
the inequality it is meant to prove, the event domain, and the error
budget. This manuscript does not identify an already working method for
that question. RSCT-E01 remains reserved rather than an executable proof
task; the conceptual investigation remains open.

\subsection{Test stronger event information jointly}

Full all-realness contains more information than
\eqref{eq:physicaljet}. Higher deflated Laguerre inequalities, additional
event jets for \(m>2\), or constraints on a complete noncentral zero
measure could supply further necessary conditions.
Their utility is not settled by the present constructions.
For each proposed addition one must derive the physical condition,
transform the carrier exactly, pay the required derivatives, and test
the enlarged set of conditions at one common point.
Separate compatibility witnesses do not prove joint compatibility.
Conversely, the one successful joint construction here does not imply
that every finite family of conditions is ineffective.

\subsection{Distinguish useful mean estimates from exceptional events}

The neighboring Green-potential program has positive results for an
actual canonical observable: at \(T=1/L_X\), its dyadic moments tend to
zero for every fixed positive power, with corresponding subcritical
exponential-moment conclusions \cite{gnp19}. These are moving-time
average statements, with no present rate or every-center conclusion.
At an all-real threshold its Green value is already zero, even at a
multiple real zero \cite{gnp20}; the value alone cannot exclude the
threshold event.

Three tools from this investigation are suitable for discussion there:
the explicit dependence on a growing calibrated prime set; the complete
local approximation and derivative-error propagation; and the
calibrated phase constructions as tests of proposed estimates.
Their original measurements and quantifiers must be retained during a
transfer. In particular, the auxiliary phase average in
Lemma~\ref{lem:complement} is not an actual-height average.

A concrete proposed question is to quantify the Green program's
fixed-cutoff residual mean
\[
 \frac1U\int_U^{2U}R_P(X,1)\dd X=\mu_P+E(U,P),
 \qquad
 \mu_P=\zeta(3/2)\prod_{p\le P}(1-p^{-3/2})-1,
\]
uniformly for a declared growing range of \(P\), ideally including a
cutoff of order \((\log U)^2\). Here \(R_P\) is the complete normalized
two-arm residual of the Green program, not \(\cD_a\).
This is an unproved quantitative question. A cutoff fixed on each
dyadic block but growing between blocks is an intermediate target;
passing to a cutoff varying with \(X\) inside the integral requires
another argument. Even a successful mean estimate leaves sparse
exceptional centers.
\repo{papers/investigations-5/09_green_potentials_newman_flow/01_riemann_siegel_evolution/brainstorm/external/01_RSCT_MANUSCRIPT_COMPARISON_AND_PROPOSED_QUESTIONS_20261112.md}{The companion comparison note}
gives the exact residual definition,
source links, and the additional obligations for a local deflation
transfer. No such transfer is claimed proved here.

\subsection{Cover the full event domain}

Even a favorable result throughout \eqref{eq:band} would cover only a
specified asymptotic regime. The original threshold problem includes
bounded centers, other values of \(T\log(a/(4\pi))\), and arbitrary
multiplicity. Known fixed-positive-time high-zero reality and simplicity
are inherited inputs, not new coverage from this investigation.
A complete argument must state the remaining domain partition and
prove all required overlaps without applying fixed-geometry estimates
to changing geometries.

\subsection{Audit inputs and make constants effective where useful}

The field provider, its larger complex charts, Gaussian exterior, and
fractional-cutoff correction need an independent specialist audit before
the present conditional proof chain is treated as independently
established. The local and matrix constants are finite, but their
numerical values and onsets have not been certified. Making them
effective is useful only for a clearly specified application with an
available arithmetic margin.
Finite computations could then support that application by checking
normalizations, evaluating complete expressions with errors, or covering
a bounded remaining region. They cannot replace the missing
all-center arithmetic argument or validate it by sampled agreement.
